% GENERATED FILE. Edit canonical lessons, then run npm run generate:books.
% canonical-source-hash: fnv1a64:942204e4c5cbc8a7
% canonical-lessons: TA-C187-cellappirani, TA-C187-kalai, TA-C187-kanru, TA-C187-onay, TA-C187-ciruttai, TA-R187-first-pass-words-from-chapters-179-182, TA-R187-second-pass-words-from-chapters-183-187

\chapter{Pet, Bull, Calf, Wolf, Leopard}
\label{ch:ta-pet-bull-calf-wolf-leopard}
\hlchaptermodality{187}

\begin{chapteropening}
\textbf{By the end of this chapter:} \emph{I can say a pet, a bull, a calf, a wolf and a leopard.}

Complete the last lesson of chapter 187: I can say a pet, a bull, a calf, a wolf and a leopard.
\end{chapteropening}

\section[\texorpdfstring{\ta{செல்லப்பிராணி} (cellappirāṇi)}{செல்லப்பிராணி (cellappirāṇi)}]{\ta{செல்லப்பிராணி} (cellappirāṇi) — a pet}

\label{lesson:TA-C187-cellappirani}



\begin{quote}
Before the new one: say the Tamil for romantic love, then the Tamil for friendship.
\end{quote}

\subsection*{You'll want to know: \ta{செல்லப்பிராணி}}
\textbf{\ta{செல்லப்பிராணி}} — \emph{cellappirāṇi} — \textquotedblleft{}a pet\textquotedblright{}.

\begin{grammarlens}[title={What you've built}]
One more word for this chapter.
\end{grammarlens}

\subsection*{Your turn}
\begin{itemize}
\raggedright
  \item \emph{Say it:} \emph{cellappirāṇi}
  \item \emph{Say it:} \emph{cellappirāṇi}, once more
  \item \emph{Say it:} say \emph{cellappirāṇi}
  \item \emph{Recall:} say the Tamil for romantic love, then the Tamil for friendship, then say \emph{cellappirāṇi} again
\end{itemize}

\begin{tcolorbox}[breakable,colback=teal!4,colframe=teal!35!black,title={Before you move on}]
What does \ta{செல்லப்பிராணி} mean? (\textquotedblleft{}A pet\textquotedblright{}.) Say it once more.
\end{tcolorbox}
\section[\texorpdfstring{\ta{காளை} (kāḷai)}{காளை (kāḷai)}]{\ta{காளை} (kāḷai) — a bull}

\label{lesson:TA-C187-kalai}



\begin{quote}
Before the new one: say the Tamil for friendship, then the Tamil for a pet.
\end{quote}

\subsection*{You'll want to know: \ta{காளை}}
\textbf{\ta{காளை}} — \emph{kāḷai} — \textquotedblleft{}a bull\textquotedblright{}.

\begin{grammarlens}[title={What you've built}]
One more word for this chapter.
\end{grammarlens}

\subsection*{Your turn}
\begin{itemize}
\raggedright
  \item \emph{Say it:} \emph{kāḷai}
  \item \emph{Say it:} \emph{kāḷai}, once more
  \item \emph{Say it:} say \emph{kāḷai}
  \item \emph{Recall:} say the Tamil for friendship, then the Tamil for a pet, then say \emph{kāḷai} again
\end{itemize}

\begin{tcolorbox}[breakable,colback=teal!4,colframe=teal!35!black,title={Before you move on}]
What does \ta{காளை} mean? (\textquotedblleft{}A bull\textquotedblright{}.) Say it once more.
\end{tcolorbox}
\section[\texorpdfstring{\ta{கன்று} (kaṉṟu)}{கன்று (kaṉṟu)}]{\ta{கன்று} (kaṉṟu) — a calf}

\label{lesson:TA-C187-kanru}



\begin{quote}
Before the new one: say the Tamil for a pet, then the Tamil for a bull.
\end{quote}

\subsection*{You'll want to know: \ta{கன்று}}
\textbf{\ta{கன்று}} — \emph{kaṉṟu} — \textquotedblleft{}a calf\textquotedblright{}.

\begin{grammarlens}[title={What you've built}]
One more word for this chapter.
\end{grammarlens}

\subsection*{Your turn}
\begin{itemize}
\raggedright
  \item \emph{Say it:} \emph{kaṉṟu}
  \item \emph{Say it:} \emph{kaṉṟu}, once more
  \item \emph{Say it:} say \emph{kaṉṟu}
  \item \emph{Recall:} say the Tamil for a pet, then the Tamil for a bull, then say \emph{kaṉṟu} again
\end{itemize}

\begin{tcolorbox}[breakable,colback=teal!4,colframe=teal!35!black,title={Before you move on}]
What does \ta{கன்று} mean? (\textquotedblleft{}A calf\textquotedblright{}.) Say it once more.
\end{tcolorbox}
\section[\texorpdfstring{\ta{ஓநாய்} (ōnāy)}{ஓநாய் (ōnāy)}]{\ta{ஓநாய்} (ōnāy) — a wolf}

\label{lesson:TA-C187-onay}



\begin{quote}
Before the new one: say the Tamil for a bull, then the Tamil for a calf.
\end{quote}

\subsection*{You'll want to know: \ta{ஓநாய்}}
\textbf{\ta{ஓநாய்}} — \emph{ōnāy} — \textquotedblleft{}a wolf\textquotedblright{}.

\begin{grammarlens}[title={What you've built}]
One more word for this chapter.
\end{grammarlens}

\subsection*{Your turn}
\begin{itemize}
\raggedright
  \item \emph{Say it:} \emph{ōnāy}
  \item \emph{Say it:} \emph{ōnāy}, once more
  \item \emph{Say it:} say \emph{ōnāy}
  \item \emph{Recall:} say the Tamil for a bull, then the Tamil for a calf, then say \emph{ōnāy} again
\end{itemize}

\begin{tcolorbox}[breakable,colback=teal!4,colframe=teal!35!black,title={Before you move on}]
What does \ta{ஓநாய்} mean? (\textquotedblleft{}A wolf\textquotedblright{}.) Say it once more.
\end{tcolorbox}
\section[\texorpdfstring{\ta{சிறுத்தை} (ciṟuttai)}{சிறுத்தை (ciṟuttai)}]{\ta{சிறுத்தை} (ciṟuttai) — a leopard}

\label{lesson:TA-C187-ciruttai}



\begin{quote}
Before the new one: say the Tamil for a calf, then the Tamil for a wolf.
\end{quote}

\subsection*{You'll want to know: \ta{சிறுத்தை}}
\textbf{\ta{சிறுத்தை}} — \emph{ciṟuttai} — \textquotedblleft{}a leopard\textquotedblright{}.

\begin{grammarlens}[title={What you've built}]
One more word for this chapter.
\end{grammarlens}

\subsection*{Your turn}
\begin{itemize}
\raggedright
  \item \emph{Say it:} \emph{ciṟuttai}
  \item \emph{Say it:} \emph{ciṟuttai}, once more
  \item \emph{Say it:} say \emph{ciṟuttai}
  \item \emph{Recall:} say the Tamil for a calf, then the Tamil for a wolf, then say \emph{ciṟuttai} again
\end{itemize}

\begin{tcolorbox}[breakable,colback=teal!4,colframe=teal!35!black,title={Before you move on}]
What does \ta{சிறுத்தை} mean? (\textquotedblleft{}A leopard\textquotedblright{}.) Say it once more.
\end{tcolorbox}
\section[(dialogue)]{Words from chapters 179-182 — a first pass}

\label{lesson:TA-R187-first-pass-words-from-chapters-179-182}



\begin{quote}
Say the words for a wolf and a leopard.
\end{quote}

\subsection*{Your turn}
Say these aloud:

\begin{itemize}
\raggedright
  \item useful, local, foreign, clear, a feeling
  \item active, naughty, stubborn, funny, interesting
  \item favourite, terrible, many, everything, thick
  \item a refrigerator, a tap, a tumbler, a pan, a lid
\end{itemize}

\begin{tcolorbox}[breakable,colback=teal!4,colframe=teal!35!black,title={Before you move on}]
Useful? (\textbf{\ta{பயனுள்ள}}, \emph{payaṉuḷḷa}.) Local? (\textbf{\ta{உள்ளூர்}}, \emph{uḷḷūr}.) Foreign? (\textbf{\ta{வெளிநாட்டு}}, \emph{veḷināṭṭu}.) Clear? (\textbf{\ta{தெளிவான}}, \emph{teḷivāṉa}.) A feeling? (\textbf{\ta{உணர்ச்சி}}, \emph{uṇarcci}.) Active? (\textbf{\ta{சுறுசுறுப்பான}}, \emph{cuṟucuṟuppāṉa}.) Naughty? (\textbf{\ta{குறும்பான}}, \emph{kuṟumpāṉa}.) Stubborn? (\textbf{\ta{பிடிவாதமான}}, \emph{piṭivātamāṉa}.) Funny? (\textbf{\ta{நகைச்சுவையான}}, \emph{nakaiccuvaiyāṉa}.) Interesting? (\textbf{\ta{சுவாரசியமான}}, \emph{cuvāraciyamāṉa}.) Favourite? (\textbf{\ta{பிடித்தமான}}, \emph{piṭittamāṉa}.) Terrible? (\textbf{\ta{மோசமான}}, \emph{mōcamāṉa}.) Many? (\textbf{\ta{பல}}, \emph{pala}.) Everything? (\textbf{\ta{எல்லாம்}}, \emph{ellām}.) Thick? (\textbf{\ta{தடிமனான}}, \emph{taṭimaṉāṉa}.) A refrigerator? (\textbf{\ta{குளிர்சாதனப்} \ta{பெட்டி}}, \emph{kuḷircātaṉap peṭṭi}.) A tap? (\textbf{\ta{குழாய்}}, \emph{kuḻāy}.) A tumbler? (\textbf{\ta{டம்ளர்}}, \emph{ṭamḷar}.) A pan? (\textbf{\ta{சட்டி}}, \emph{caṭṭi}.) A lid? (\textbf{\ta{மூடி}}, \emph{mūṭi}.)
\end{tcolorbox}
\section[(dialogue)]{Words from chapters 183-187 — a second pass}

\label{lesson:TA-R187-second-pass-words-from-chapters-183-187}



\begin{quote}
Say the words for pongal and upma.
\end{quote}

\subsection*{Your turn}
Say these aloud:

\begin{itemize}
\raggedright
  \item pongal, upma, sambar, a chutney, a dry stir-fried vegetable dish
  \item ice, breakfast, lunch, dinner, a notebook
  \item a queue, a number, a text message, news; a message, an email
  \item a father-in-law, a boyfriend, a girlfriend, romantic love, friendship
  \item a pet, a bull, a calf, a wolf, a leopard
\end{itemize}

\begin{tcolorbox}[breakable,colback=teal!4,colframe=teal!35!black,title={Before you move on}]
Pongal? (\textbf{\ta{பொங்கல்}}, \emph{poṅkal}.) Upma? (\textbf{\ta{உப்புமா}}, \emph{uppumā}.) Sambar? (\textbf{\ta{சாம்பார்}}, \emph{cāmpār}.) A chutney? (\textbf{\ta{சட்னி}}, \emph{caṭṉi}.) A dry stir-fried vegetable dish? (\textbf{\ta{பொரியல்}}, \emph{poriyal}.) Ice? (\textbf{\ta{பனிக்கட்டி}}, \emph{paṉikkaṭṭi}.) Breakfast? (\textbf{\ta{காலை} \ta{உணவு}}, \emph{kālai uṇavu}.) Lunch? (\textbf{\ta{மதிய} \ta{உணவு}}, \emph{matiya uṇavu}.) Dinner? (\textbf{\ta{இரவு} \ta{உணவு}}, \emph{iravu uṇavu}.) A notebook? (\textbf{\ta{குறிப்பேடு}}, \emph{kuṟippēṭu}.) A queue? (\textbf{\ta{வரிசை}}, \emph{varicai}.) A number? (\textbf{\ta{எண்}}, \emph{eṇ}.) A text message? (\textbf{\ta{குறுஞ்செய்தி}}, \emph{kuṟuñceyti}.) News; a message? (\textbf{\ta{செய்தி}}, \emph{ceyti}.) An email? (\textbf{\ta{மின்னஞ்சல்}}, \emph{miṉṉañcal}.) A father-in-law? (\textbf{\ta{மாமனார்}}, \emph{māmaṉār}.) A boyfriend? (\textbf{\ta{காதலன்}}, \emph{kātalaṉ}.) A girlfriend? (\textbf{\ta{காதலி}}, \emph{kātali}.) Romantic love? (\textbf{\ta{காதல்}}, \emph{kātal}.) Friendship? (\textbf{\ta{நட்பு}}, \emph{naṭpu}.) A pet? (\textbf{\ta{செல்லப்பிராணி}}, \emph{cellappirāṇi}.) A bull? (\textbf{\ta{காளை}}, \emph{kāḷai}.) A calf? (\textbf{\ta{கன்று}}, \emph{kaṉṟu}.) A wolf? (\textbf{\ta{ஓநாய்}}, \emph{ōnāy}.) A leopard? (\textbf{\ta{சிறுத்தை}}, \emph{ciṟuttai}.)
\end{tcolorbox}
